\section{Power system stabilizer --- \texttt{pss2a}}
\label{sec:pss2a}

\subsection{Purpose of a power system stabilizer}

A power system stabilizer (PSS) is a supplementary control loop on the excitation
system of a synchronous generator. Its only task is to damp the electromechanical
oscillations of the rotor against the rest of the grid; it contributes nothing to
the steady-state voltage or power setpoints and is inactive when the machine runs
at constant speed.

The need for it arises from the voltage regulator itself. A modern static or
brushless excitation system with a high-gain, fast-acting automatic voltage
regulator (AVR) such as the \texttt{esac8b} of Section~\ref{sec:esac8b} holds the
terminal voltage tightly and thereby raises the synchronising torque of the machine,
which improves transient stability after a fault. The same fast field action,
however, produces a component of electrical torque that lags the rotor speed
deviation, which reduces --- and in weakly connected systems can eliminate --- the
natural damping of the rotor swing. The result is a growing or poorly damped
oscillation of rotor angle, power and voltage in the range of roughly 0.1 to 2\,Hz:
local plant modes near 1--2\,Hz, in which a unit or a plant swings against the
nearby system, and inter-area modes near 0.1--0.8\,Hz, in which groups of machines
in different regions swing against each other~[3].

The PSS restores this damping. It measures a signal that reflects the rotor speed
deviation and, after phase compensation, adds a small modulating signal $V_S$ to the
voltage reference of the AVR. The phase compensation is tuned so that the resulting
change in electrical torque is in phase with the speed deviation; a torque component
in phase with speed is by definition a damping torque. The output is limited (here to
$\pm0.1$\,pu) so that the stabilizer can only ever nudge the terminal voltage and
cannot override the voltage regulator during a large disturbance.

Early stabilizers used speed alone as the input, which made them sensitive to
torsional shaft modes, or electrical power alone, which caused an unwanted
terminal-voltage excursion whenever the turbine output was ramped. The
\emph{integral-of-accelerating-power} design standardised as IEEE~421.5
type~PSS2A~[4] combines both: speed and electrical power are used to synthesise the
integral of the accelerating power, $\int (P_{mech} - P_{elec})\,dt$, which is
proportional to the speed deviation but free of the torsional content and
insensitive to slow changes in mechanical power. This is the design applied to all
nine Cape Station units.

For the Fervo units the stabilizer serves two purposes. It provides the damping that
the fast \texttt{esac8b} regulator would otherwise remove, and it satisfies the WECC
requirement that synchronous generators equipped with suitable excitation systems be
fitted with a stabilizer and have it in service, which is a condition of model
acceptance for the interconnection studies this data set feeds.

\subsection{Model implementation}

\texttt{pss2a} is the IEEE~421.5 dual-input stabilizer. Input~1 is rotor speed
deviation (\texttt{j1} $= 1$) and input~2 is electrical power (\texttt{j2} $= 3$),
both measured locally. The two signals are washed out ($T_{w1} = T_{w2} = T_{w3} =
2$\,s, $T_{w4} = 0$ so the fourth washout is bypassed), the power signal is integrated
through $K_{S2}/(1 + s \cdot T_7)$ with $T_7 = 2$\,s and $K_{S2} = 0.501$, and the two
are combined and passed through the ramp-tracking filter with $T_8 = 0.5$\,s,
$T_9 = 0.1$\,s, $N = 1$, $M = 5$ to synthesise the integral of accelerating power.
The result is amplified by $K_{S1} = 5$ and phase-compensated by two identical
lead-lag stages $(1 + 0.25 s)/(1 + 0.02 s)$, then limited to $\pm0.1$\,pu before it
enters the exciter as $V_S$. Figure~\ref{fig:pss2a} shows the structure.

The value $K_{S2} = 0.501$ is consistent with the usual tuning rule
$K_{S2} = T_7 / (2 \cdot H)$ evaluated with $H = 2.0$\,s, which does not match either
machine inertia ($H = 1.49$ or $1.50$\,s would give $K_{S2} = 0.671$ or $0.667$). This
is worth confirming with the supplier of the dynamic data.

\begin{figure}[!ht]
\centering
\resizebox{\linewidth}{!}{\begin{tikzpicture}[x=1cm,y=1cm]
  % input 1 row
  \node[blk,minimum width=18mm] (w1) at (0.8,0) {$\dfrac{s \cdot T_{w1}}{1 + s \cdot T_{w1}}$};
  \node[blk,minimum width=18mm] (w2) at (3.0,0) {$\dfrac{s \cdot T_{w2}}{1 + s \cdot T_{w2}}$};
  \node[blk,minimum width=16mm] (t6) at (5.2,0) {$\dfrac{1}{1 + s \cdot T_6}$};
  \node[sum] (sa) at (6.8,0) {};
  \node[blk,minimum width=28mm] (rt) at (8.8,0) {$\left[\dfrac{1 + s \cdot T_8}{(1 + s \cdot T_9)^M}\right]^N$};
  \node[sum] (sb) at (10.8,0) {};
  \node[blk,minimum width=10mm] (ks1) at (12.0,0) {$K_{S1}$};
  \node[blk,minimum width=18mm] (l1) at (13.7,0) {$\dfrac{1 + s \cdot T_1}{1 + s \cdot T_2}$};
  \node[blk,minimum width=18mm] (l2) at (15.9,0) {$\dfrac{1 + s \cdot T_3}{1 + s \cdot T_4}$};
  \node[lim,minimum width=9mm] (lim) at (17.6,0) {};
  \draw[thick] ($(lim.center)+(-0.25,-0.25)$) -- ($(lim.center)+(-0.1,-0.25)$) -- ($(lim.center)+(0.1,0.25)$) -- ($(lim.center)+(0.25,0.25)$);
  % input 2 row
  \node[blk,minimum width=18mm] (w3) at (0.8,-2.0) {$\dfrac{s \cdot T_{w3}}{1 + s \cdot T_{w3}}$};
  \node[blk,minimum width=18mm] (w4) at (3.0,-2.0) {$\dfrac{s \cdot T_{w4}}{1 + s \cdot T_{w4}}$};
  \node[blk,minimum width=16mm] (t7) at (5.2,-2.0) {$\dfrac{K_{S2}}{1 + s \cdot T_7}$};
  \node[blk,minimum width=10mm] (ks3) at (6.8,-2.0) {$K_{S3}$};
  \node[blk,minimum width=10mm] (ks4) at (10.8,-2.0) {$K_{S4}$};
  % signals
  \draw[sig] (-0.9,0) node[left,lbl]{$\Delta\omega$} -- (w1);
  \draw[sig] (w1) -- (w2);
  \draw[sig] (w2) -- (t6);
  \draw[sig] (t6) -- (sa);
  \draw[sig] (sa) -- (rt);
  \draw[sig] (rt) -- (sb);
  \node[lbl] at ($(sb)+(-0.22,-0.22)$) {$-$};
  \draw[sig] (sb) -- (ks1);
  \draw[sig] (ks1) -- (l1);
  \draw[sig] (l1) -- (l2);
  \draw[sig] (l2) -- (lim);
  \node[lbl,above=0.5mm of lim] {$V_{STMAX}$};
  \node[lbl,below=0.5mm of lim] {$V_{STMIN}$};
  \draw[sig] (lim) -- (18.8,0) node[right,lbl]{$V_S$};
  \draw[sig] (-0.9,-2.0) node[left,lbl]{$P_{elec}$} -- (w3);
  \draw[sig] (w3) -- (w4);
  \draw[sig] (w4) -- (t7);
  \draw[sig] (t7) -- (ks3);
  \draw[sig] (ks3.north) -- (sa);
  \fill (6.15,-2.0) circle (1pt);
  \draw[sig] (6.15,-2.0) -- (6.15,-3.0) -- (10.8,-3.0) -- (ks4.south);
  \draw[sig] (ks4.north) -- (sb);
\end{tikzpicture}
}
\caption{\texttt{pss2a} block diagram. $T_{w4} = 0$ bypasses the fourth washout; $K_{S3} = K_{S4} = 1$.}
\label{fig:pss2a}
\end{figure}

\begin{table}[!ht]
\centering
\begin{tabularx}{\linewidth}{@{}l X l r@{}}
\toprule
\textbf{Parameter} & \textbf{Description} & \multicolumn{1}{c}{\textbf{Unit}} & \multicolumn{1}{c}{\textbf{Value}} \\
\midrule
\texttt{j1} & Input signal 1 code (1 = rotor speed deviation) & -- & 1 \\
\texttt{k1} & Remote bus for input 1 (0 = local) & -- & 0 \\
\texttt{j2} & Input signal 2 code (3 = electrical power) & -- & 3 \\
\texttt{k2} & Remote bus for input 2 (0 = local) & -- & 0 \\
\texttt{tw1} & Washout time constant $T_{w1}$ & s & 2 \\
\texttt{tw2} & Washout time constant $T_{w2}$ & s & 2 \\
\texttt{tw3} & Washout time constant $T_{w3}$ & s & 2 \\
\texttt{tw4} & Washout time constant $T_{w4}$ (0 = bypassed) & s & 0 \\
\texttt{t6} & Input 1 transducer time constant $T_6$ & s & 0 \\
\texttt{t7} & Input 2 transducer time constant $T_7$ & s & 2 \\
\texttt{ks2} & Input 2 gain $K_{S2}$ ($= T_7 / (2 \cdot H)$) & pu & 0.501 \\
\texttt{ks3} & Signal-matching gain $K_{S3}$ & pu & 1 \\
\texttt{ks4} & Input 2 gain $K_{S4}$ & pu & 1 \\
\texttt{t8} & Ramp-tracking filter lead time constant $T_8$ & s & 0.5 \\
\texttt{t9} & Ramp-tracking filter lag time constant $T_9$ & s & 0.1 \\
\texttt{n} & Ramp-tracking filter order $N$ & -- & 1 \\
\texttt{m} & Ramp-tracking filter order $M$ & -- & 5 \\
\texttt{ks1} & Stabilizer gain $K_{S1}$ & pu & 5 \\
\texttt{t1} & Lead time constant $T_1$ & s & 0.25 \\
\texttt{t2} & Lag time constant $T_2$ & s & 0.02 \\
\texttt{t3} & Lead time constant $T_3$ & s & 0.25 \\
\texttt{t4} & Lag time constant $T_4$ & s & 0.02 \\
\texttt{vstmax} & Stabilizer output upper limit $V_{STMAX}$ & pu & 0.1 \\
\texttt{vstmin} & Stabilizer output lower limit $V_{STMIN}$ & pu & -0.1 \\
\texttt{a} & Optional second-order filter coefficient $a$ & -- & 0 \\
\texttt{ta} & Optional filter time constant $T_a$ & s & 0 \\
\texttt{tb} & Optional filter time constant $T_b$ & s & 0 \\
\bottomrule
\end{tabularx}

\caption{\texttt{pss2a} parameters (common to all nine units).}
\label{tab:pss2a}
\end{table}
